\begin{examplebox}
    \textbf{\textcolor{CExample}{例1（基础：外点直接求切点弦）}}

    \textbf{题目：} 椭圆 $\frac{x^2}{4}+\frac{y^2}{3}=1$，点 $P(4,3)$。求从 $P$ 引两条切线的切点弦方程。
    \begin{description}[style=nextline, leftmargin=2.8em, labelsep=0.5em, itemsep=0.35em]
        \item[\textbf{第1步（验证是外点）：}] $S=\frac{4^2}{4}+\frac{3^2}{3}=7>1$，所以有两个不同切点。
        \item[\textbf{第2步（写方程）：}] 将 $a^2=4,b^2=3,x_0=4,y_0=3$ 代入，得到
        \[\frac{4x}{4}+\frac{3y}{3}=1\ \Longrightarrow\ \boxed{x+y=1}.\]
        \item[\textbf{第3步（自查）：}] 把外点代入左边得 $4+3=7\ne1$；符合``切点弦不经过外点''。
    \end{description}

    \medskip
    \textbf{\textcolor{CExample}{例2（常见分母：规范化简）}}

    \textbf{题目：} 椭圆 $\frac{x^2}{9}+\frac{y^2}{4}=1$，外点 $P(6,2)$，求切点弦方程。
    \begin{description}[style=nextline, leftmargin=2.8em, labelsep=0.5em, itemsep=0.35em]
        \item[\textbf{第1步：}] $S=\frac{36}{9}+\frac4{4}=5>1$。
        \item[\textbf{第2步：}] $\frac{6x}{9}+\frac{2y}{4}=1$，化简为 $\frac{2x}{3}+\frac y2=1$。
        \item[\textbf{第3步：}] 两端乘 $6$，得到 $\boxed{4x+3y=6}$。
    \end{description}

    \medskip
    \textbf{\textcolor{CExample}{例3（逆向：已知切点弦反求外点）}}

    \textbf{题目：} 椭圆 $\frac{x^2}{4}+y^2=1$ 的一条切点弦为 $2x+y=1$，求相应的外点 $P$。
    \begin{description}[style=nextline, leftmargin=2.8em, labelsep=0.5em, itemsep=0.35em]
        \item[\textbf{第1步：}] 写出 $\frac{x_0}{4}x+y_0y=1$。
        \item[\textbf{第2步：}] 与 $2x+y=1$ 比较系数，得 $\frac{x_0}{4}=2,\ y_0=1$，故 $P(8,1)$。
        \item[\textbf{第3步：}] 验证 $S=\frac{64}{4}+1=17>1$，符合外点条件。
    \end{description}
    \textbf{变式提醒：} 若已知切点弦为 $Ax+By=C$（$C\ne0$），先改写为
    $\frac ACx+\frac BCy=1$，再得
    $\boxed{x_0=\frac{a^2A}{C},\ y_0=\frac{b^2B}{C}}$。

    \medskip
    \textbf{\textcolor{CExample}{例4（轴上外点：竖直的切点弦）}}

\begin{center}
\begin{tikzpicture}[x=1.1cm,y=1.1cm,font=\small,
  every node/.style={inner sep=3pt},line cap=round]

\draw[->,gray!65] (-3.4,0)--(6.6,0) node[right] {$x$};
\draw[->,gray!65] (0,-2.4)--(0,2.6) node[above] {$y$};
\draw[CStatement,very thick] (0,0) ellipse [x radius=3,y radius=2];
\draw[orange!85!black,thick] (1.5,1.73205)--(6,0)--(1.5,-1.73205);
\draw[CTrap,very thick] (1.5,-2.1)--(1.5,2.1);
\fill (1.5,1.73205) circle (1.5pt) node[above right,yshift=3pt] {$A(\frac32,\sqrt3)$};
\fill (1.5,-1.73205) circle (1.5pt) node[below right,yshift=-3pt] {$B(\frac32,-\sqrt3)$};
\fill (1.5,0) circle (1.5pt) node[below right] {$M(\frac32,0)$};
\fill (6,0) circle (1.5pt) node[above] {$P(6,0)$};
\node[below left] at (0,0) {$O$};

\end{tikzpicture}

{\small 例4：$P(6,0)$ 的切点弦是竖直线 $x=\frac32$。}
\end{center}

    \textbf{题目：} 椭圆 $\frac{x^2}{9}+\frac{y^2}{4}=1$，点 $P(6,0)$，求切点弦方程及其中点。
    \begin{description}[style=nextline, leftmargin=2.8em, labelsep=0.5em, itemsep=0.35em]
        \item[\textbf{第1步：}] $S=\frac{36}{9}+0=4>1$。
        \item[\textbf{第2步：}] $\frac{6x}{9}+0=1$，故 $\boxed{x=\frac32}$。
        \item[\textbf{第3步：}] $M=(\frac 64,\frac04)=(\frac32,0)$。
        \item[\textbf{第4步（检验有两个交点）：}] 将 $x=\frac32$ 代入椭圆，
        $\frac{y^2}{4}=1-\frac14=\frac34$，故切点为
        $\bigl(\frac32,\sqrt3\bigr)$、$\bigl(\frac32,-\sqrt3\bigr)$。
    \end{description}

    \medskip
    \textbf{\textcolor{CExample}{例5（最小反例：内点不能套成切点弦）}}

    \textbf{题目：} 对椭圆 $\frac{x^2}{4}+y^2=1$，若 $P(1,0)$，能否将 $\frac{x}{4}=1$ 称为切点弦？

    \textbf{解：} $S=\frac14<1$，$P$ 在椭圆内，不能作实切线。
    形式代换虽给出 $x=4$，但 $x=4$ 与椭圆没有交点，故\textbf{不是切点弦}。

    \medskip
    \textbf{\textcolor{CExample}{高价值变式（先自行尝试）}}

    椭圆 $\frac{x^2}{16}+\frac{y^2}{9}=1$ 的切点弦所在直线为 $3x+2y=6$。
    求外点 $P$ 与切点弦中点 $M$。

    \textbf{答案与思路：} 先化成 $\frac x2+\frac y3=1$，所以
    $P(8,3)$；$S=\frac{64}{16}+\frac99=5$，故
    $\boxed{M(\frac85,\frac35)}$。检验 $P$ 在椭圆外，因为 $S=5>1$。
\end{examplebox}
